% Header include for TESTS_AND_DIAGNOSTICS pandoc -> xelatex build.
% Goal: prevent overlapping text in wide tables.

\usepackage{array}
\usepackage{etoolbox}

% Tighter column padding so long cells don't overflow
\setlength{\tabcolsep}{3pt}
\renewcommand{\arraystretch}{1.15}

% Shrink table fonts so wide rows fit the page
\AtBeginEnvironment{longtable}{\footnotesize}
\AtBeginEnvironment{tabular}{\footnotesize}
\AtBeginEnvironment{tabularx}{\footnotesize}

% Allow URLs / inline code to break across lines
\usepackage{xurl}
\renewcommand{\UrlBreaks}{\do\/\do\-\do\_\do\.\do\=}

% pandoc's `inline code` becomes \texttt{...} — long paths overflow
% table columns. The Lua filter .pandoc_break_paths.lua rewrites
% these as \BreakablePath{...} with \discretionary{}{}{} break
% opportunities already inserted after / _ - . — so we only need to
% set tt-family and allow sloppy line breaking.
\newcommand{\BreakablePath}[1]{{\ttfamily\sloppy #1}}

% --- Symbol fallback for glyphs missing from STIX Two Text -----------------
% STIX Two Text lacks several common math/symbol glyphs in the OT:scri
% feature set we render with (→, ≤, ≥, ∈, ≈, ↔, ★). Map each one to its
% counterpart in STIX Two Math, which has the full set.
\usepackage{newunicodechar}
\newfontfamily\stixfallback{STIX Two Math}[Scale=MatchLowercase]
\newcommand{\fb}[1]{{\stixfallback #1}}

\newunicodechar{→}{\fb{→}}
\newunicodechar{←}{\fb{←}}
\newunicodechar{↔}{\fb{↔}}
\newunicodechar{≤}{\fb{≤}}
\newunicodechar{≥}{\fb{≥}}
\newunicodechar{∈}{\fb{∈}}
\newunicodechar{≈}{\fb{≈}}
\newunicodechar{★}{\fb{★}}
\newunicodechar{⊕}{\fb{⊕}}
\newunicodechar{σ}{\fb{σ}}
\newunicodechar{Σ}{\fb{Σ}}
\newunicodechar{Δ}{\fb{Δ}}
\newunicodechar{α}{\fb{α}}
\newunicodechar{β}{\fb{β}}
\newunicodechar{γ}{\fb{γ}}
\newunicodechar{δ}{\fb{δ}}
\newunicodechar{ε}{\fb{ε}}
\newunicodechar{η}{\fb{η}}
\newunicodechar{λ}{\fb{λ}}
\newunicodechar{μ}{\fb{μ}}
\newunicodechar{π}{\fb{π}}
\newunicodechar{ρ}{\fb{ρ}}
\newunicodechar{τ}{\fb{τ}}
\newunicodechar{φ}{\fb{φ}}
\newunicodechar{χ}{\fb{χ}}
\newunicodechar{²}{\fb{²}}
\newunicodechar{³}{\fb{³}}
\newunicodechar{ω}{\fb{ω}}
\newunicodechar{Ω}{\fb{Ω}}
\newunicodechar{∝}{\fb{∝}}
\newunicodechar{√}{\fb{√}}
\newunicodechar{∞}{\fb{∞}}
\newunicodechar{±}{\fb{±}}
\newunicodechar{×}{\fb{×}}
\newunicodechar{÷}{\fb{÷}}
\newunicodechar{−}{\fb{−}}
\newunicodechar{☉}{\fb{☉}}
\newunicodechar{✓}{\fb{✓}}
\newunicodechar{⚠}{\fb{⚠}}
